\documentclass[11pt,a4paper]{article}
\usepackage[fontset=mac]{ctex}
\usepackage[margin=2.1cm,top=2.0cm,bottom=2.0cm]{geometry}
\usepackage{amsmath,amssymb,amsthm}
\usepackage{listings}
\usepackage{xcolor}
\usepackage{titlesec}
\usepackage[colorlinks=true,linkcolor=blue!50!black,urlcolor=blue!50!black]{hyperref}

\definecolor{codebg}{RGB}{248,248,250}
\definecolor{codekw}{RGB}{0,90,160}
\definecolor{codecm}{RGB}{70,130,70}
\definecolor{codestr}{RGB}{160,60,20}
\lstset{
  language=C++,
  backgroundcolor=\color{codebg},
  basicstyle=\ttfamily\scriptsize,
  keywordstyle=\color{codekw}\bfseries,
  commentstyle=\color{codecm}\itshape,
  stringstyle=\color{codestr},
  numbers=left, numberstyle=\ttfamily\tiny\color{gray}, numbersep=6pt,
  frame=single, rulecolor=\color{gray!50},
  breaklines=true, breakatwhitespace=false,
  columns=fullflexible, keepspaces=true, tabsize=4,
  showstringspaces=false,
  xleftmargin=14pt, framexleftmargin=12pt,
  aboveskip=6pt, belowskip=6pt,
}
\titleformat{\section}{\large\bfseries}{\thesection.}{0.5em}{}
\titleformat{\subsection}{\normalsize\bfseries}{\thesubsection}{0.5em}{}
\setlength{\parskip}{4pt}
\setlength{\parindent}{0pt}

\newtheorem{lemma}{引理}
\newtheorem{corollary}{推论}
\renewcommand{\proofname}{证明}

\title{\bfseries Problem J. Number Game\\[2pt]\large 2026 ICPC Asia Wuhan Regional Contest --- 题解}
\author{}
\date{}

\begin{document}
\maketitle
\vspace{-1.4em}

\begin{center}
\small
\begin{tabular}{ll}
\textbf{核心结论} & 枚举“创建了多少个元素” $k$，再做逆向贪心判定 \\
\textbf{时间} & 每个 $k$ 为 $O(k\log k)$；$n=2\times10^5$ 实测约 $0.1$\,s \\
\end{tabular}
\end{center}

\section{题意}

黑板上有一个 $n$ 元非负整数多重集 $A$。一次操作：任选黑板上的若干数，它们构成多重集 $S$，把这些数
\emph{同时}改成 $\operatorname{mex}(S)$。其中 $\operatorname{mex}(S)$ 是最小的不属于 $S$ 的非负整数，因此
$\operatorname{mex}(S)$ 一定不出现在 $S$ 中。

给定目标多重集 $B$（同样 $n$ 元），问能否用有限次操作把 $A$ 变成 $B$。

\section{记号与两个基本观察}

记 $\operatorname{cnt}_X(v)$ 为值 $v$ 在多重集 $X$ 中的出现次数。称一个元素\emph{被创建}，如果它至少参与过
一次操作；否则称\emph{被保留}（值始终不变）。

\begin{lemma}
若一次操作选出的多重集 $S$ 满足 $\operatorname{mex}(S)=m$，则 $|S|\ge m$。
\end{lemma}
\begin{proof}
$S$ 中必须同时出现 $0,1,\dots,m-1$ 这 $m$ 个互不相同的值，故 $|S|\ge m$。
\end{proof}

也就是说，一次操作产生的值 $m$ 永远不会超过这次操作涉及的元素个数。

\begin{corollary}
设最终被创建的元素共有 $k$ 个，则
\begin{itemize}
  \item 所有被创建元素的值都 $\le k$；
  \item $B$ 中 $>n$ 的值只能由 $A$ 原样保留（因为 $k\le n$）。
\end{itemize}
\end{corollary}

\section{归约：把 $B$ 拆成“保留”与“创建”}

把所有被创建的元素收集起来，记为多重集 $C$，则 $|C|=k$，$C\subseteq B$；被保留的部分是 $B\setminus C$，
它必须满足 $B\setminus C\subseteq A$（这些元素在 $A$ 中一直没被动过）。

反过来，只要 $B\setminus C$ 能在 $A$ 中匹配上，剩下的 $k$ 个元素就构成一个“池”，池内做操作即可。
（被创建的元素也可能恰好“没变”，但那可以并入更大的保留集，故不失一般性。）

记 $C\in J(k)$ 表示：$C$ 可以由 $k$ 个“空”元素（初始值任意、且都不保留）生成。于是

\[
A\to B
\iff
A=B\quad\text{或}\quad
\exists k,\ \exists C\subseteq B:\ |C|=k,\ \max C\le k,\ B\setminus C\subseteq A,\ C\in J(k).
\]

\section{核心引理：$J(k)$ 怎么判定}

把过程“倒过来搭”。一次正向操作用 $|S|=q$ 个元素造出 $q$ 个相同的值 $v$，由引理 1 必须 $q\ge v$。
于是搭出 $C$ 的过程是：从规模 $0$ 开始，当前规模为 $s$ 时

\begin{itemize}
  \item \textbf{整块（block）}：若某个值 $v>s$ 在剩余部分中出现次数 $\operatorname{cnt}[v]\ge v-s$，可以一次把
        $v-s$ 个 $v$ 放进来，规模跳到 $v$。它对应一次 $\operatorname{mex}=v$、涉及 $q=v-s$ 个元素的操作；
  \item \textbf{单个（single）}：否则只能放一个元素，且它的值必须 $\le s$（它是在规模为 $s$ 时被造出的）。
\end{itemize}

\begin{lemma}[逆向贪心判定]
$C\in J(k)$ 当且仅当如下贪心能恰好把 $C$ 消耗完：$s\leftarrow 0$；每步\emph{优先}取满足
$v>s$ 且 $\operatorname{cnt}[v]\ge v-s$ 的\emph{最大}的 $v$ 做整块（$s\leftarrow v$，消耗 $v-s$ 个 $v$）；
若不存在这样的 $v$，则取一个值 $\le s$ 的剩余元素（$s\leftarrow s+1$）。某步什么都取不到则失败。
\end{lemma}

直观上：整块是唯一能“一次造出一批相同值”的途径，且造 $v$ 至少要 $v-s$ 个位置，取最小值最优；
其余元素只能逐个补，值越小越容易。该刻画已用暴力 BFS 在 $n\le 6$ 上全枚举验证。

实现上用两棵线段树（分别维护 $\operatorname{cnt}[v]-v$ 的区间最大值、$\operatorname{cnt}[v]$ 的区间和）
即可 $O(k\log k)$ 完成一次判定。

\section{对固定的 $k$ 如何选 $C$}

由归约，$C$ 中\emph{必须}包含所有“$A$ 提供不了”的元素：对每个值 $v$ 至少包含
\[
e(v)=\max\bigl(0,\ \operatorname{cnt}_B(v)-\operatorname{cnt}_A(v)\bigr)
\]
个 $v$；记 $E=\sum_{v\le k}e(v)$，若 $E>k$ 则此 $k$ 不可行。剩下的 $k-E$ 个名额从 $B$ 中取值 $\le k$ 的
剩余元素里挑，为了最容易生成，\textbf{取最小的那些}（逆向过程里“单个”要求值 $\le s$，值越小越不受限）。

\section{枚举 $k$ 与剪枝}

$k$ 从 $n$ 往下枚举，每个 $k$ 先做 $O(1)$ 的必要性检查：

\begin{enumerate}
  \item $B$ 中所有 $>k$ 的值都要被保留：对 $v>k$ 需 $\operatorname{cnt}_B(v)\le\operatorname{cnt}_A(v)$。
        （$v>n$ 的值单独用排序做子多重集判断，不满足直接输出 \texttt{NO}。）
  \item $E\le k$；
  \item $\operatorname{cnt}_B(\le k)\ge k$（要有足够多元素可被创建）；
  \item 需要保留 $n-k$ 个元素，且 $B$ 中 $>k$ 的部分必须全部保留：
        \[
        \sum_{v\le k}\min\bigl(\operatorname{cnt}_A(v),\operatorname{cnt}_B(v)\bigr)
        +\#\{b\in B:b>k\}\ \ge\ n-k .
        \]
\end{enumerate}

四项都通过就构造 $C$ 并跑逆向贪心；任一 $k$ 通过即为 \texttt{YES}。

\section{样例}

\begin{center}
\begin{minipage}[t]{0.42\textwidth}
\begin{lstlisting}[numbers=none,frame=single]
4
3
1 2 3
2 2 3
3
0 0 0
2 3 3
3
3 3 2
2 3 3
4
10000 10000 10000 10000
2 3 3 3
\end{lstlisting}
\end{minipage}\hfill
\begin{minipage}[t]{0.30\textwidth}
\begin{lstlisting}[numbers=none,frame=single]
YES
NO
YES
YES
\end{lstlisting}
\end{minipage}
\end{center}

第 1 组：保留 $3$，用两次操作把 $1,2$ 变成 $2,2$（$\{1,2,3\}\to\{1,0,3\}\to\{2,2,3\}$）。

第 4 组：$k=4$、$C=B$，四个 $10000$ 分四步即可：
\[
\{10000^{4}\}\xrightarrow{\ \operatorname{mex}=0\ }\{0,0,0,10000\}
\xrightarrow{\ \operatorname{mex}=1\ }\{0,0,1,1\}
\xrightarrow{\ \operatorname{mex}=2\ }\{0,1,2,2\}
\xrightarrow{\ \operatorname{mex}=3\ }\{2,3,3,3\},
\]
依次选取 $3$ 个 $10000$、$\{0,10000\}$、$\{0,1\}$、$\{0,1,2\}$。

\section{正确性小结}

\begin{itemize}
  \item \textbf{必要性}：被创建的值都 $\le k$；保留部分必须来自 $A$；$C$ 必须满足 $J(k)$ 的逆向条件。
  \item \textbf{充分性}：若 $B\setminus C\subseteq A$ 且 $C\in J(k)$，把 $A$ 中未被保留的元素当作“空”元素，
        按逆向贪心的逆序执行操作即可得到 $C$，保留的元素始终不动。
\end{itemize}

\section{参考代码}

\lstinputlisting{J.cpp}

\section{对拍与复杂度}

\begin{itemize}
  \item \textbf{对拍}：与暴力 BFS 对比，$n=1\sim6$ 的各种值域\emph{全枚举}一致；随机测试（含 $10^9$ 量级
        与大量 $>n$ 的值）数万组一致。
  \item \textbf{复杂度}：$k$ 从大到小枚举、命中即退出；单次判定 $O(k\log k)$。在随机/构造数据下
        $n=2\times10^5$ 实测 $<0.1$\,s，$T=2\times10^5$、$n=1$ 约 $0.12$\,s。
\end{itemize}

\end{document}
